\section{Smooth reservoirs and the physical boundary law}
\label{sec:boundary-smooth}

\subsection{What depends on the reservoir entropy}

Let a reservoir induce a log likelihood $h_N(E)$ relative to a specified
canonical law. For reservoir entropy $s_{B,N}=S_{B,N}/k_B$,
$h_N(E)=\beta_NE+s_{B,N}(\mathcal E_N-E)$ up to a constant.
Where temperature and positive heat capacity are defined by this
entropy convention,
\begin{equation}
 -h_N''(E)=-s_{B,N}''(\mathcal E_N-E)
 =\frac{\beta_B(\mathcal E_N-E)^2}{C_B(\mathcal E_N-E)/k_B}.
 \label{eq:smooth-capacity}
\end{equation}
The proofs thus concern entropy curvature, with a heat-capacity
interpretation when these thermodynamic derivatives exist.

\begin{proposition}[Curvature necessity]
\label{prop:smooth-necessity}
Suppose a positive decomposition \eqref{eq:positive-decomposition}
has two phases with positive limiting weights, centers
$e_{-,N}<e_{+,N}$ separated by $\Delta_N\asymp N$, and conditional
limits $(E-e_{i,N})/\sqrt N\Rightarrow\Normal(0,v_i)$ with $v_i>0$.
Let $Q_N$ have likelihood proportional to $e^{h_N}$, with $h_N=-\infty$
outside an interval. On its interior assume $h_N$ is twice continuously
differentiable and $-h_N''\ge k\kappa_N$ between the phase centers,
for some $k>0$ and a positive sequence $\kappa_N>0$.
Then $\TV(P_N,Q_N)\to0$ requires $\kappa_NN^{3/2}\to0$.
\end{proposition}
\begin{proof}
TV convergence gives $h_N(E)-\log Z_N\to0$ in canonical probability
on the feasible set. Choose good energies with standardized positions
in $[1/2,1]$ and $[3/2,2]$ above the lower center, and $[-1,-1/2]$
below the upper center. Each interval has positive limiting probability.
The selected energies $x_N<y_N<z_N$ are feasible, their log likelihoods
differ from $\log Z_N$ by $o(1)$, and
$y_N-x_N\asymp\sqrt N$, $z_N-y_N\asymp N$.
Their convex hull is feasible and lies between the centers.
Concavity of $h_N(E)+k\kappa_NE^2/2$ gives
\[
 h_N(y_N)-\frac{z_N-y_N}{z_N-x_N}h_N(x_N)
 -\frac{y_N-x_N}{z_N-x_N}h_N(z_N)
 \ge \frac{k\kappa_N}{2}(y_N-x_N)(z_N-y_N).
\]
The left side tends to zero. This is the curvature version of
Theorem~\ref{thm:two-scale-necessity}'s good-point chord argument.
\end{proof}

Under the density hypotheses of Section~\ref{subsec:density-tails},
Proposition~\ref{prop:density-sufficiency} supplies the converse when
$h_N(e_{-,N})=h_N(e_{+,N})=0$, $h_N$ is globally concave on its
feasible interval, and its curvature satisfies
\eqref{eq:concave-curvature-upper}. A family with curvature comparable
to $\kappa_N$ therefore has the sharp $\kappa_NN^{3/2}\to0$ criterion.
This statement assumes that endpoint calibration exists; an arbitrary
entropy need not realize it by varying only the composite energy.
For exact separated Gaussian phases no interphase density hypothesis
is needed: the chord and curvature upper bound give vanishing
standardized endpoint slopes, and the tangent bound is uniformly
integrable by Gaussian exponential moments. This gives the analogous
criterion $\kappa_N\Delta_Ns_N\to0$ for a gap $\Delta_N$ and Gaussian
fluctuation scale $s_N$.

\subsection{Finite tilts at capacity proportional to \texorpdfstring{$N^{3/2}$}{N to the power 3/2}}

Here $\beta_N\to\beta>0$. Use the positive decomposition
\eqref{eq:positive-decomposition}, with $w_{i,N}\to w_i>0$,
$w_-+w_+=1$, and $\Delta_N/N\to\ell>0$.
Write $\mathcal E_N^{\mathrm{sec}}$ for
\eqref{eq:secant-calibration}. A Gaussian measure
$G_i=\Normal(0,v_i)$ may have $v_i=0$, interpreted as a point mass.

\begin{theorem}[Physical boundary and finite calibration corrections]
\label{thm:physical-boundary}
Suppose $c_N/N^{3/2}\to\gamma\in(0,\infty)$ and
\begin{equation}
 X_{i,N}:=(E-e_{i,N})/\sqrt N\Rightarrow G_i,\qquad
 \sup_{N\ge N_0}\E_{P_{i,N}}e^{\eta|X_{i,N}|}<\infty
 \quad\text{for some }N_0\text{ and }\eta>b:=\frac{\beta^2\ell}{2\gamma}.
 \label{eq:boundary-moment}
\end{equation}
Assume also
\begin{equation}
 \delta_N\exp\!\left[\max_E h_N^{\mathrm{sec}}(E)\right]\longrightarrow0.
 \label{eq:boundary-exception}
\end{equation}
For $\mathcal E_N=\mathcal E_N^{\mathrm{sec}}+d_N\sqrt N$, $d_N\to d$,
put $b_-=b$, $b_+=-b$, $D=\beta^2\ell d/\gamma$, and
\begin{equation}
 A_-=0,\quad A_+=D,\qquad
 Z(d)=\sum_iw_i e^{A_i+b^2v_i/2},\qquad
 r_i(d)=\frac{w_i e^{A_i+b^2v_i/2}}{Z(d)}.
 \label{eq:boundary-phase-weights}
\end{equation}
The reservoir phase probabilities, defined on the positive auxiliary
decomposition, tend to $r_i(d)$; their standardized conditional
fluctuations converge weakly to $G_i^b:=\Normal(b_iv_i,v_i)$.
Macroscopic energy separation makes these limiting probabilities
observable from energy alone. On the full microscopic state space,
\begin{equation}
 \lim_N\TV(P_N,Q_N)
 =\frac12\sum_iw_i\int_{\R}
 \left|\frac{e^{A_i+b_i z}}{Z(d)}-1\right|\dd G_i(z).
 \label{eq:physical-boundary-tv}
\end{equation}
\end{theorem}
\begin{proof}
Normalize $h_N$ at $e_{-,N}$. Uniformly on bounded standardized windows,
the exact secant formula gives
\begin{equation}
 h_N(e_{i,N}+\sqrt Nz)=A_i+b_i z+o(1).
 \label{eq:physical-local-tilt}
\end{equation}
Indeed, with $q_N=\beta_N\Delta_N/c_N\to0$, the secant endpoint
derivatives are $\beta_Nq_N/2+O(q_N^2)$ and
$-\beta_Nq_N/2+O(q_N^2)$. Their products with $\sqrt N$ tend to
$b$ and $-b$, and $N\sup|h_N''|=O(N/c_N)\to0$ on these windows.
Changing $\mathcal E_N$ by $d_N\sqrt N$ changes the upper-minus-lower
endpoint value by
\[
 d_N\sqrt N\,c_N
 \left[(\mathcal E_N^{\mathrm{sec}}-e_{+,N})^{-1}
 -(\mathcal E_N^{\mathrm{sec}}-e_{-,N})^{-1}\right]+o(1)
 \longrightarrow D,
\]
and changes each standardized endpoint slope by $o(1)$.

Concavity gives
$h_N(E)\le h_N(e_{i,N})+h_N'(e_{i,N})(E-e_{i,N})$
on the feasible set. Hence
$e^{h_N(E)}\le C\exp[(b+\epsilon)|X_{i,N}|]$ for small fixed
$\epsilon>0$ and large $N$. Choose $p>1$ with $p(b+\epsilon)<\eta$.
Equation~\eqref{eq:boundary-moment} gives uniform integrability, also
for the absolute difference in \eqref{eq:physical-boundary-tv}.
Cutoff points have zero weight and cause no difficulty.
The maximum of $h_N$ changes from the secant maximum by $O(1)$ for
bounded $d_N$. Indeed, the maximum is at
$E=\mathcal E_N-c_N/\beta_N$ and, when $h_N(e_{-,N})=0$, the derivative
of its value with respect to $\mathcal E_N$ is
$\beta_N-c_N/(\mathcal E_N-e_{-,N})=O(N^{-1/2})$ in this range.
Thus \eqref{eq:boundary-exception} removes the exceptional contribution.
Weak convergence, \eqref{eq:physical-local-tilt}, and uniform
integrability now give $Z_N\to Z(d)$ and the Gaussian tilts.
Finally use the exact identity
$\TV(P_N,Q_N)=\frac12\E_{P_N}|e^{h_N}/Z_N-1|$.
This does not infer an integral limit from weak convergence alone.
\end{proof}

The secant prescription $d=0$ equalizes likelihoods at the centers,
but its limiting integrated phase weights differ from $w_i$ if the
variances differ. If both variances equal $v$, the weights remain
unchanged and
\begin{equation}
 \lim_N\TV(P_N,Q_N)=2\Phi(b\sqrt v/2)-1,
 \label{eq:physical-balanced-boundary}
\end{equation}
where $\Phi$ is the standard normal distribution function.

\begin{corollary}[Compensating the phase probabilities]
\label{cor:physical-compensation}
Under these hypotheses, the calibration
\begin{equation}
 \mathcal E_N=\mathcal E_N^{\mathrm{sec}}
 +\frac{\beta^2\ell(v_--v_+)}{8\gamma}\sqrt N+o(\sqrt N)
 \label{eq:physical-compensation}
\end{equation}
restores the original phase weights asymptotically, and
\begin{equation}
 \lim_N\TV(P_N,Q_N)
 =\sum_iw_i\{2\Phi(b\sqrt{v_i}/2)-1\}.
 \label{eq:compensated-tv}
\end{equation}
\end{corollary}
\begin{proof}
The correction makes $D=b^2(v_--v_+)/2$, so the integrated factors
in \eqref{eq:boundary-phase-weights} are equal. Within each phase
the remaining change is the Gaussian translation by $b_iv_i$, with
TV $2\Phi(b\sqrt{v_i}/2)-1$. This is zero when $v_i=0$.
\end{proof}
Matching phase populations thus leaves a finite distortion of phase
fluctuations. The correction does not assert exact finite-$N$ matching.

\subsection{Optimization over every composite energy}

For finite signed measures let $\|\cdot\|_{\mathrm{var}}$ be the total
variation norm. Define the continuous convex function
\begin{equation}
 F(r)=\frac12\left\{
 \|w_-G_- -rG_-^b\|_{\mathrm{var}}
 +\|w_+G_+ -(1-r)G_+^b\|_{\mathrm{var}}\right\},
 \qquad 0\le r\le1.
 \label{eq:physical-boundary-objective}
\end{equation}

\begin{theorem}[Optimized physical boundary]
\label{thm:physical-boundary-optimum}
Under Theorem~\ref{thm:physical-boundary}'s hypotheses,
\begin{equation}
 \lim_N\inf_{\mathcal E:\,Z_N(\mathcal E)>0}\TV(P_N,Q_{N,\mathcal E})
 =\min\left\{\min_{0\le r\le1}F(r),\,w_-,\,w_+\right\}.
 \label{eq:physical-optimal-boundary}
\end{equation}
For $w_-=w_+=1/2$ and $v_-=v_+=v$, this is
$\min\{2\Phi(b\sqrt v/2)-1,1/2\}$.
\end{theorem}
\begin{proof}
As $d$ varies through $\R$, \eqref{eq:boundary-phase-weights} realizes
every strictly positive pair of phase weights. The boundary theorem
therefore realizes $F(r)$ for every $r\in(0,1)$; continuity supplies
the upper bound by its infimum on $[0,1]$.
The additional upper bounds $w_-$ and $w_+$ follow by centering the
reservoir maximum at either phase, as in
\eqref{eq:centered-calibration}. Its weight lies in $[0,1]$, tends to
one in the selected phase since $N/c_N\to0$, and tends to zero in the
other since $\Delta_N^2/c_N\to\infty$. Positive decomposition and
macroscopic separation give TV tending to the discarded phase weight.

For the lower bound consider a subsequence of arbitrary calibrations
with TV limit strictly below $\min(w_-,w_+)$. Let $R_N$ be their
normalized likelihood. The overlap in phase $i$ is
$O_{i,N}=w_{i,N}\E_{P_{i,N}}\min(R_N,1)$.
Since $\TV(P_N,Q_N)=\E_{P_N}(1-R_N)_+$, each phase loss is at most TV;
hence $O_{i,N}\ge w_{i,N}-\TV(P_N,Q_N)$. Both overlaps stay positive.
Tightness, the overlap bound from $R_N<\epsilon$, and
$P_N(R_N>M)\le1/M$ allow selection of feasible points
$x_N=e_{-,N}+O(\sqrt N)$ and $y_N=e_{+,N}+O(\sqrt N)$ with
$R_N(x_N),R_N(y_N)\in[\epsilon,M]$, for fixed positive $\epsilon,M$.
Set $a_N=\log R_N(y_N)-\log R_N(x_N)=O(1)$ and
$D_N=y_N-x_N\asymp N$. The exact likelihood solves to
\begin{equation}
 \mathcal E_N=x_N+
 \frac{D_N}{1-\exp[-(\beta_ND_N-a_N)/c_N]}.
 \label{eq:boundary-overlap-calibration}
\end{equation}
A bounded $a_N$ changes this by $O(c_N/D_N)=O(\sqrt N)$ from
the secant calibration for $x_N,y_N$. Uniformly on these windows
that calibration is
\[
 \frac{c_N}{\beta_N}+\frac{x_N+y_N}{2}
 +\frac{\beta_N(y_N-x_N)^2}{12c_N}
 +O\!\left(\frac{(y_N-x_N)^4}{c_N^3}\right).
\]
Moving its endpoints by $O(\sqrt N)$ changes it by $O(\sqrt N)$.
Thus $(\mathcal E_N-\mathcal E_N^{\mathrm{sec}})/\sqrt N$ is bounded.
Pass to a convergent subsequence and apply
Theorem~\ref{thm:physical-boundary}; the limit is $F(r)\ge\min F$.
Sequences with limiting error at least $\min(w_-,w_+)$ already obey
the required bound. Apply this argument to calibrations within $1/N$
of the infimum; no existence of minimizers is assumed.
The upper and lower bounds coincide, so the limit in
\eqref{eq:physical-optimal-boundary} exists and has the stated value.
In the balanced equal-variance case reflection sends $r$ to $1-r$,
so convexity minimizes $F$ at $1/2$, giving the normal formula.
\end{proof}

\noindent
The optimization is a scalar calculation. Set $r_-=r$, $r_+=1-r$,
and $d_i=b\sqrt{v_i}$. For $d_i>0$ the contribution of phase $i$
to the positive part of reservoir minus reference,
$r_iG_i^b-w_iG_i$, is
\begin{equation}
 r_i\Phi\!\left(\frac{\log(r_i/w_i)}{d_i}+\frac{d_i}{2}\right)
 -w_i\Phi\!\left(\frac{\log(r_i/w_i)}{d_i}-\frac{d_i}{2}\right).
 \label{eq:physical-optimum-cdf}
\end{equation}
Their sum is $F(r)$ because total positive and negative masses agree.
Use zero for $r_i=0$, by continuity, and $(r_i-w_i)_+$ when $d_i=0$.
Equation~\eqref{eq:physical-optimum-cdf} follows by integrating on the
half-line where the tilted Gaussian density exceeds the weighted
original density. If both $d_i>0$, differentiation gives
\[
 F'(r)=
 \Phi\!\left(\frac{\log(r/w_-)}{d_-}+\frac{d_-}{2}\right)
 -\Phi\!\left(\frac{\log((1-r)/w_+)}{d_+}+\frac{d_+}{2}\right).
\]
This is strictly increasing from a negative to a positive value.
Its unique zero, equivalently equality of the two CDF arguments,
determines the unique minimizer. For equal positive variances it is
$r=w_-$, even with unequal target weights. Thus the optimized limit
then simplifies to
$\min\{2\Phi(b\sqrt v/2)-1,w_-,w_+\}$.

\subsection{An exact physical gain and its interfacial scale}

At a fraction $\theta$ of the gap, the endpoint-normalized physical
secant likelihood has the exact form
\begin{equation}
 h_N(e_{-,N}+\theta\Delta_N)
 =c_N\left[\theta q_N+
 \log(1-\theta+\theta e^{-q_N})\right],
 \qquad q_N=\frac{\beta_N\Delta_N}{c_N}.
 \label{eq:physical-interior-exact}
\end{equation}
For $q_N\to0$, uniformly on $0\le\theta\le1$,
\begin{align}
 h_N(e_{-,N}+\theta\Delta_N)
 &=\frac{\beta_N^2\Delta_N^2}{2c_N}\theta(1-\theta)
 \notag\\
 &\quad+\frac{\beta_N^3\Delta_N^3}{6c_N^2}
       \theta(1-\theta)(2\theta-1)
 +O(\Delta_N^4/c_N^3).
 \label{eq:physical-interior-expansion}
\end{align}
This follows by expanding the logarithm of the Bernoulli transform
in \eqref{eq:physical-interior-exact}; bounded $\theta$ gives a uniform
remainder. The maximum lies inside the two centers because the
endpoint slopes have opposite signs. Equation~\eqref{eq:physical-interior-expansion}
therefore proves \eqref{eq:boundary-maximal-gain}.
At $c_N\asymp N^{3/2}$ the cubic term can be of order one, so the
quadratic approximation does not determine all subleading weights.

An interface with surface cost of order $L^{d-1}=N^{(d-1)/d}$ competes
with a maximal bath gain of order $N^2/c_N$.
Their balance occurs at $c_N\asymp N^{1+1/d}$.
For $d>2$ this is below the $N^{3/2}$ fluctuation threshold.
For $d=2$ they coincide. This comparison explains both the exponent
$\alpha=1/2$ in the density criterion and the restricted two-dimensional
boundary range \eqref{eq:two-dimensional-boundary-range}.
It is a scale comparison; a surface bound alone does not identify
a favored droplet or slab.

\subsection{Microscopic applications of the boundary law}
\label{subsec:boundary-applications}

The boundary law holds for the mean-field model of
Theorem~\ref{thm:mf-bath}, for every $\gamma>0$ and kinetic parameter
$a\ge0$. The occupation bound in Lemma~\ref{lem:mf-exponential}
is Gaussian on the $\sqrt N$ scale, so it bounds
$\E e^{t|X_{i,N}|}$ for every fixed $t$ and sufficiently large $N$.
The independent Gamma contribution has the same property by its
exact transform. There is no exceptional component.
The variances are \eqref{eq:mf-variances}; for pure spins $v_+=0$,
which is allowed here. A continuous kinetic energy is unnecessary
for the boundary theorem or its optimization.

For the short-range model of Theorem~\ref{thm:sr-bath}, the boundary
law and its optimization hold for every $\gamma>0$ in fixed spatial
dimension $d\ge3$, and for all sufficiently large fixed $\gamma$ in
$d=2$. The additional argument is needed because a small exponential
moment does not justify every finite tilt.
The cutoff-free restricted-partition window in
Lemma~\ref{lem:sr-restricted-derivatives} has width $\eta_0/L$.
When $d>2$, every fixed multiple of $N^{-1/2}=L^{-d/2}$ eventually
lies in that window. The second derivative bound
\eqref{eq:sr-restricted-second} and center estimate
\eqref{eq:sr-restricted-center} therefore bound exponential moments
of the standardized bond count at every fixed parameter.
The inverse-temperature to bond-fugacity change has bounded derivatives
and preserves this conclusion. The conditional binomial estimate and
Cauchy--Schwarz argument in Lemma~\ref{lem:sr-energy-transfer}
transfer these bounds to the actual spin energy.
The contour-phase weak limits agree with
\eqref{eq:sr-conditional-clt}: the already proved small exponential
moment bounds the conditional probability of the opposite macroscopic
midpoint phase by $Ce^{-c\sqrt N}$, and the exceptional contour
probability vanishes. Positive limiting phase weights then identify
the conditional laws.

The exceptional probability is at most $Ce^{-\tau L^{d-1}}$.
The exact secant likelihood satisfies
\begin{equation}
 \frac{\max_Eh_N^{\mathrm{sec}}(E)}{\sqrt N}
 \longrightarrow\frac{\beta^2\ell^2}{8\gamma}.
 \label{eq:boundary-maximal-gain}
\end{equation}
For $d>2$ the surface exponent dominates $\sqrt N$, giving
\eqref{eq:boundary-exception} for every $\gamma$.
For $d=2$, let $\eta_*>0$ be a valid spin-phase moment coefficient
and $\tau>0$ a valid exceptional surface constant. It suffices that
\begin{equation}
 \frac{\beta^2\ell}{2\gamma}<\eta_*,
 \qquad \frac{\beta^2\ell^2}{8\gamma}<\tau.
 \label{eq:two-dimensional-boundary-range}
\end{equation}
These hold for all sufficiently large $\gamma$. No assertion about
all $\gamma>0$ or small-capacity microscopic morphologies in two
dimensions follows from these estimates.

An alternative sufficient input is the continuous density envelope
\eqref{eq:interior-density-envelope} with $\alpha>1/2$, together with
\eqref{eq:local-density-assumption}. These also give the boundary
law and finite corrections for every fixed $\gamma$.
On the interior put $x=\min_i|E-e_{i,N}|$. Concavity gives
$h_N^{\mathrm{sec}}(E)\le Cx/\sqrt N$ at this boundary.
With a finite calibration correction, the endpoint values are bounded
but need not both vanish, so the bound instead reads
$h_N(E)\le C_0+C_1x/\sqrt N$ with constants independent of $N$.
For $x\ge R\sqrt N$ the ratio of the linear term to $x^2/N$ is at
most $C_1/R$, and its ratio to $x^\alpha$ is at most
$C_1N^{1/2-\alpha}$. Choose $R$ large and then $N$ large to absorb
this term in the envelope. The offset contributes only the bounded
prefactor $e^{C_0}$ to the tail integral and does not affect its
vanishing limit as $R\to\infty$.
Exterior to the centers the secant weight is at most one.
For finite corrections, concavity and the endpoint slope signs instead
bound exterior weights by a constant, which suffices.
Weighted tail control and local density convergence give the same
integral formula. At $\alpha=1/2$ the second ratio need not vanish.
The capillarity model in Appendix~\ref{sec:capillarity-diagnostics}
explains this distinction.
